\section{Rate transfer at matched clocks}
\label{sec:rate-transfer}

We keep the saturated homogeneous component $p_\eps$ selected in
Section~\ref{sec:characteristic-plane}, with least aspect $m>0$.
Its information and time profiles are $f,S$, and its matched-plane
potential is $\Phi=f\circ\Pi_m$ from
\eqref{eq:matched-plane-profile}.  The row data and their admissible
depths are those of \eqref{eq:trace-variables} and
\eqref{eq:trace-domain}.  Write $x,y,w,z$ for their plane trace rates
and $D$ for the label rate.  When $m=1$, the middle datum is
$U=(Y,W)$ and its single rate is $q=a_U$.  The matched component
defining $f,S,\Phi$ is fixed before their trace points.  Later
homogeneous selections concern the derived block laws and do not
redefine these profiles.

\begin{proposition}[Matched-clock rate transfer]
\label{prop:rate-transfer}
For almost every point $(i,u)$ of the characteristic plane with $i>0$
and $0<u<1$, put $s=S'(u)$.  There is a constant $\sigma>0$,
depending only on the matched process and $m$, for which $0<s\le1$
and the following assertions hold.
\begin{enumerate}
\item If $m\ne1$, each of the chains $x\longrightarrow y$ and
$w\longrightarrow z$ has the rule
\[
 d>0\quad\Longrightarrow\quad e\ge\min(1,d+s),
\]
where $d,e$ are its source and target rates.  Each of
$x\longrightarrow w$ and $y\longrightarrow z$ has the same rule
with $\sigma$ in place of $s$.
\item If $m=1$, then
\[
 x>0\Longrightarrow q\ge s,\qquad
 x>0,\ z<1\Longrightarrow z\ge x+s,\qquad
 z<1\Longrightarrow q\le2-s.
\]
\end{enumerate}
The bounds include $0\le x,y,w,z\le1$, or $0\le q\le2$ in
the tied case, supplied by the coordinate capacities.
\end{proposition}

The proof first extracts a finite point--parameter law from an
affine entropy block.  That law will carry the parameter
nonconcentration, coordinate support counts, and product comparisons
used by the projection tests.  Distinct clocks give the four
rate transfers directly.  At tied clocks, the joint rate $q$ is
supplemented by directions between parameter samples and two
conditional projection tests.

\subsection{The selected profiles and the parameter children}

We first verify what the already selected sequence $p_\eps$ supplies
at the depths needed here.  Its initial homogeneous list contains all
finite tuples from a countable dense family of admissible configurations,
including the indicated conditional probabilities given $T$.
For two nearby configurations, give both descriptions the finer label
clock.  The box comparisons in Section~\ref{sec:characteristic-plane}
give bounded ambiguity for changing the data frame at a fixed admissible
resolution.  A depth change of size $a$ introduces at most
$\eps^{-Ca+o(1)}$ possibilities, also conditional on $T$, on each fixed
compact depth range.  Applying the conditional-surprise counting bound
from Lemma~\ref{lem:profile-diagonal} to both descriptions extends the
homogeneous exponents from the dense family to every fixed admissible
configuration.  The same argument applies to finite tuples.  It does
not require independent data refinements to remain admissible at the
coarse label clock.

Write $\mathfrak h(A)$ for the limiting normalized entropy of a listed
observation.  For a finite pair $A,B$, the selected homogeneous
probabilities give, in probability,
\begin{equation}
 -\logeps{p_\eps(B\mid A)}
 =\mathfrak h(A,B)-\mathfrak h(A)+o(1).
 \label{eq:rate-ratio}
\end{equation}
The same statement holds for listed probabilities conditional on $T$;
the probability of $T$ itself is never binned.  After a trace point,
a positive block length $h$, and a finite list of real configurations
have been chosen, take one event $\mathcal E_\eps$ on which all their original
conditional surprises differ from their limiting exponents by at most
$\omega_\eps=o(h)$.  Its probability tends to one.  On this event,
ratios of the joint and parent masses give uniform upper conditional
masses and upper retained support counts.  The support count follows
by dividing the upper mass of a parent by the lower mass of each
retained joint value.

This finite real-depth event is imposed before any restriction which
uses its bounds.  If a later law satisfies
$\mathsf Q\le Kp_\eps(\,\cdot\mid\mathcal E_\eps)$, with
$\log K=o(h\ell)$ and $\ell=\log(1/\eps)$, its conditional entropy
of each listed child given its listed parent differs from the listed exponent by
at most $\omega_\eps+\log(K/p_\eps(\mathcal E_\eps))/\ell$, by
conditional relative entropy and the support count.  The parent
likelihood and child-support thresholds in the proof of
Lemma~\ref{lem:homogenization} give the same absolute error bound for
realized surprise, with $\log(1/\eta)/\ell$ added, outside
$\mathsf Q$-mass at most $2\eta$; choose $\eta\downarrow0$ with
$\log(1/\eta)=o(h\ell)$.  The dense-list
profiles remain the same under this high-probability restriction by
Lemma~\ref{lem:profile-diagonal}.  This verifies the finite inheritance
for the component selected in Section~\ref{sec:characteristic-plane};
it does not select a new $f$ or assert uniform bins at all real depths.

For this component the prediction identities and the positive constant
obtained from its least-aspect argument are
\begin{equation}
 \Ent(R_{mv}\mid T,C_v)=0,\qquad
 \Ent(R_{m(v+a)}\mid T,C_v)\ge\kappa a
 \quad(0<v<v+a<1).
 \label{eq:rate-prediction}
\end{equation}
These are limiting normalized statements for this same $p_\eps$.

Fix a regular plane point $(i,u)$ and a trace block length $h>0$.
Write $E$ for the coarse observation at $\Pi_m(i,u)$ and put
$\rho=\eps^h$.  Given the coarse label, $E$ is a function of
$(T,C_u,R_{mu})$ up to bounded grid choices, which we record when
needed.  Choose $K>\max(1,1/m)$, fixed throughout this construction,
and define the normalized parameter increments
\[
 c=\eps^{-u}(C_{u+Kh}-C_u),\qquad
 r=-\eps^{-mu}(R_{m(u+Kh)}-R_{mu}).
\]
The full finite label $\theta$ records the advanced label and the
bounded choices just mentioned.  Before all tests, fix one half-open
dyadic tree on a fixed bounded box containing the normalized pair
$(c,r)$.  For every real $0\le d\le1$, let $\Theta_d$ be its atom at
level
\[
 n_\rho(d)=\left\lceil d\log_2(1/\rho)\right\rceil.
\]
These atoms are exactly nested for all real depths: $\Theta_a$ is
determined by $\Theta_b$ when $a\le b$, and each $\Theta_a$ atom has
at most
\[
 4^{n_\rho(b)-n_\rho(a)}\le4\rho^{-2(b-a)}
\]
$\Theta_b$ children.  The time prefix records $C_{u+dh}$ and the
normal prefix records $R_{mu+dh}$ up to bounded ambiguity over $E$.
For other finite observations whose nesting holds only up to bounded
choices, we append the parent to the child when needed.  Here $h$ and
the finite prefix list are fixed before $\eps\to0$.

The chain rule gives
\begin{align}
 \Ent(C_{u+dh}\mid E)
 &\ge\Ent(C_{u+dh}\mid T,E)\notag\\
 &\ge S(u+dh)-S(u)-\Ent(R_{mu}\mid T,C_u)-o(1),
 \label{eq:rate-time-child}\\
 \Ent(R_{mu+dh}\mid E)
 &\ge\Ent(R_{mu+dh}\mid T,C_u,R_{mu})-o(1)\notag\\
 &\ge(\kappa/m)dh-o(1).
 \label{eq:rate-normal-child}
\end{align}
For the first inequality, the data in $E$ are determined by the
particle and coarse labels, and adjoining $R_{mu}$ costs at most its
conditional entropy.  The second uses
\eqref{eq:rate-prediction} at time increment $dh/m$.
Prediction also gives
\[
 \kappa a\le\Ent(R_{m(u+a)}\mid T,C_u)
       \le S(u+a)-S(u)+o(1).
\]
Consequently $s=S'(u)\ge\kappa>0$ at the regular points used below.
Fix any $0<\sigma<\min(1,\kappa/m)$.

At tied clocks there is also a bound between any two tested parameter
prefixes.  If $m=1$ and $0\le a<b\le1$, then conditional on $E$ the
symbol $\Theta_a$ is equivalent to
$(C_{u+ah},R_{u+ah})$ up to the recorded bounded choices.  Therefore
\begin{align}
 \Ent(\Theta_b\mid E,\Theta_a)
 &\ge\Ent(C_{u+bh}\mid T,E,\Theta_a)\notag\\
 &\ge S(u+bh)-S(u+ah)
       -\Ent(R_{u+ah}\mid T,C_{u+ah})-o(1)\notag\\
 &=S(u+bh)-S(u+ah)-o(1).
 \label{eq:rate-joint-child-entropy}
\end{align}
After $T,C_{u+ah}$ have been given, the extra conditioning in
$E,\Theta_a$ is a function of $R_{u+ah}$ and the bounded choices.
This is the stated entropy cost and does not use an entropy bound for $T$.

The homogeneous joint symbols $(E,\Theta_a)$ and $(E,\Theta_b)$
turn this last entropy inequality into a bound on their probability
ratio.  On the common real-depth event,
\begin{equation}
 p_\eps(\Theta_b\mid E,\Theta_a)
 \le\eps^{S(u+bh)-S(u+ah)-o(1)}
 \quad(m=1)
 \label{eq:rate-joint-child-mass}
\end{equation}
at every retained child.  The time and normal inequalities give the
corresponding marginal prefix bounds by the same ratio calculation.
Lemma~\ref{lem:typical-conditioning}, with the coarse observation
as $C$ and each listed parent containing $C$, selects base cells
of adequate relative retention without an inverse original cell mass.
We will keep the event as an unnormalized allowed set inside the
chosen base cell until both product comparisons have been imposed.

\subsection{One finite affine-block law}

Center the common pre-label data in the boxes of $E$, dilate them
to bounded coordinates, and round the resulting point $P$ at
precision $\rho$.  For a full label with normalized increment
$(c,r)$ define its canonical post-coordinate symbols by applying
the shear to this same rounded point and then taking the indicated
grid bins:
\begin{equation}
 X_\theta=X,\qquad Y_\theta=Y+cX,\qquad
 W_\theta=W+rX,\qquad
 Z_\theta=Z+cW+rY+crX.
 \label{eq:rate-shear}
\end{equation}
The same origin is transformed at every label, so the shears compose
by adding $(c,r)$ and the map from frame $0$ to frame $1$ uses
$\theta_1-\theta_0$.  Applying the shear before or after rounding
changes each real output by $O(\rho)$.  The canonical bins and the
physically recorded bins consequently determine one another with
bounded conditional ambiguity.  We use the canonical symbols in
the finite law below: each is a genuine function of $(P,\theta)$.
Their prefixes and joint numerator and denominator symbols are
included in the finite uniformization list.  Bounded ambiguity
transfers the entropy identities to these symbols; their own
uniform probability bins, rather than a pointwise assertion from
that ambiguity alone, will supply the retained mass bounds.

Let $\mathcal I=\{X,Y,W,Z\}$ when $m\ne1$, and
$\mathcal I=\{X,U,Z\}$ when $m=1$, with $U=(Y,W)$.
Write $Q_{j,\theta}$ for the canonical coordinate in group
$j\in\mathcal I$, and $(Q_{j,\theta})_d$ for its prefix at precision
$\rho^d$.  Put $(a_X,a_Y,a_W,a_Z)=(x,y,w,z)$ in the distinct-clock
case and $(a_X,a_U,a_Z)=(x,q,z)$ in the tied case.
Choose the fixed $K$ large enough that raising any
individual data depth by at most $h$ is admissible after the label
advances by $Kh$.  For any fixed finite list of $t_j\in[0,1]$, the
affine trace gives
\begin{equation}
 \Ent\bigl(((Q_{j,\theta})_{t_j})_{j\in\mathcal I}
          \mid E,\theta\bigr)
 =h\sum_{j\in\mathcal I} a_jt_j+o(h).
 \label{eq:rate-post-affine}
\end{equation}
Indeed, at the advanced label the unrefined pre- and post-data
determine one another to bounded ambiguity by
\eqref{eq:trace-domain}.  The difference between the information
with the selected refinements and this common base information is
their conditional entropy, because the particle and label determine
the refinement.  The affine trace gives the displayed increment,
and the canonical rounding costs only bounded ambiguity.

Raising all data depths by $h$ is also admissible at the original
label and is preserved under the label-conditioned shear.  The same
calculation gives
\begin{equation}
 \Ent(P\mid E)=h\sum_j a_j+o(h),\qquad
 \Ent(P\mid E,\theta)=h\sum_j a_j+o(h),\qquad
 \Mut(P;\theta\mid E)=o(h).
 \label{eq:rate-pre-independence}
\end{equation}
In particular, the conditional coordinate multiinformation obtained
by subtracting the joint entropy in \eqref{eq:rate-post-affine}
from the sum of its single-coordinate entropies is $o(h)$.
Both information quantities are nonnegative averages over the base
observation $E$.

We now form one sequence of finite blocks.  This order is needed
because the trace error divided by a fixed $h$ need not vanish as
only $\eps$ tends to zero.  Choose $h_n\downarrow0$ on the trace
subsequence for this fixed profile and plane point, so that its
affine error divided by $h_n$ tends to zero uniformly on the bounded
cone of displacements just used.  Differentiability of $S$ gives
the analogous convergence for its time increments.  At stage $n$,
include the first $n$ tuples from a countable list of rational block
prefixes, their canonical joint numerator and denominator symbols,
and the finitely many information differences above.  These are
real original depths fixed by the already chosen $h_n$.

Choose $\eps_n$ within the converging sequence $p_\eps$ so that all
finite probability, prediction, information, and box-ambiguity
errors are at most $\eta_n h_n$, where $\eta_n\downarrow0$, and
$h_n\log(1/\eps_n)\ge n$.  Put
\[
 \rho_n=\eps_n^{h_n},\qquad L_n=\log(1/\rho_n).
\]
For this extraction, let $\mathcal H_n$ contain only the tested
prefix conditions whose truth, after $E$ is fixed, is determined
by the canonical pair $(P,\theta)$: the parameter prefixes and
canonical coordinate numerator and denominator symbols.  The earlier
bins conditional on $T$ need not be pair-measurable.  They are used
to justify the entropy estimates, but are not included in
$\mathcal H_n$.
The two nonnegative information quantities and the failure
probability of the common prefix event can be selected
simultaneously on the base cells by Markov's inequality.  Choose
deterministic $d_n,\zeta_n,\omega_n\downarrow0$ so that, on every
selected base cell, the two natural-log quantities are at most
$d_nL_n$, the prefix event fails with probability at most $\zeta_n$,
and its normalized surprise errors are at most $\omega_n$ in block
units.  The finite lists and thresholds are chosen before the
precision; increasing them slowly gives these common envelopes.
Select one such cell $E=e_n$ and let $p_n^0$ be the original
conditional law of the canonical pair $(P,\theta)$ on that cell,
before normalizing the prefix event $\mathcal H_n$.  Thus
$p_n^0(\mathcal H_n)\ge1-\zeta_n$.  The cell may have small original
mass: the selected bounds are bounds for its conditional law and
contain no inverse mass of $E=e_n$.

For clarity, the two comparisons to be imposed on this law have
different spaces and factors.  Put
\[
 \mu_n^0=(p_n^0)_P,\qquad
 \lambda_n^0=(p_n^0)_\theta,\qquad
 \xi_{j,n}^{0,\theta}=(p_n^0)_{Q_{j,\theta}\mid\theta}.
\]
The first comparison is with
$\mu_n^0(P)\lambda_n^0(\theta)$.  The second is with
\[
 \lambda_n^0(\theta)
       \prod_{j\in\mathcal I}\xi_{j,n}^{0,\theta}(Q_{j,\theta});
\]
it uses the coordinate factors conditional on that label.
The divergences from these two probability laws are respectively
the pair information and the averaged conditional coordinate
multiinformation, each bounded by $d_nL_n$.  We do not assume a
uniform multiinformation estimate on individual labels.


\begin{lemma}[A finite core for a point--parameter law and its coordinates]
\label{lem:affine-block-core}
Fix an integer $k\ge1$.  Let $p^0$ be a probability law on a finite pair space
$\mathcal P\times\Theta$.  Write $\mu^0=p^0_P$ and
$\lambda^0=p^0_\theta$.  Let
$Q_j=Q_j(P,\theta)$, $1\le j\le k$, be finite coordinate
observations, and put
$\xi_j^\theta=p^0_{Q_j\mid\theta}$ on positive label fibers, choosing
an arbitrary probability law on a zero-mass label fiber.
Define the two nonnegative natural-log quantities
\begin{align*}
 D_{\mathrm{pair}}
 &=\KL(p^0_{P,\theta}\Vert\mu^0\otimes\lambda^0),\\
 D_{\mathrm{coord}}
 &=\KL\left(p^0_{\theta,Q_1,\ldots,Q_k}
       \middle\Vert
       \lambda^0(\theta)\prod_{j=1}^k\xi_j^\theta(Q_j)\right)\\
 &=\sum_{j=1}^k\Entn^{p^0}(Q_j\mid\theta)
       -\Entn^{p^0}(Q_1,\ldots,Q_k\mid\theta).
\end{align*}
The second quantity is averaged over the labels; no individual label
is assumed to have small conditional multiinformation.

Let $\mathcal F$ be a finite list of observations on the pair space
which contains $P$, $\theta$, and every $(\theta,Q_j)$.
It may also contain any finite list of prefix and joint observations.
Write $N=|\mathcal F|$, counting occurrences rather than values.
Let $\mathcal H$ be an event of $p^0$-mass at least $1-\zeta$, where
$\zeta\ge0$.
For $v>0$, choose
\[
 0<\underline M
 \le1-\zeta-\frac{D_{\mathrm{pair}}+D_{\mathrm{coord}}+2/e}{v},
 \qquad
 \delta=\frac{\underline M}{2N}.
\]
Then there is a restriction $r\le\boldsymbol1_{\mathcal H}p^0$,
of mass $M\ge\underline M/2$, such that, for every listed observation
$F$ and every value with $r_F(f)>0$,
\begin{equation}
 \delta p^0_F(f)\le r_F(f)\le p^0_F(f).
 \label{eq:block-core-floors}
\end{equation}
For the normalized law $\pi=r/M$, write
$\mu=\pi_P$ and $\lambda=\pi_\theta$.  Its pair law and its
coordinate law on every positive label fiber satisfy
\begin{align}
 \pi(P,\theta)
 &\le e^v M\delta^{-2}\mu(P)\lambda(\theta),
 \label{eq:block-core-pair}\\
 \pi(Q_1,\ldots,Q_k\mid\theta)
 &\le e^v\delta^{-k}
          \prod_{j=1}^k\pi(Q_j\mid\theta).
 \label{eq:block-core-coordinate}
\end{align}
All factors on the right of these two inequalities are marginals
of this same law $\pi$.

If $A,B$ are listed observations and $A$ is determined by $B$, then
for every retained child value $b$, with parent $a=A(b)$,
\begin{equation}
 \delta p^0(B=b\mid A=a)
 \le \pi(B=b\mid A=a)
 \le\delta^{-1}p^0(B=b\mid A=a).
 \label{eq:block-core-prefix-ratio}
\end{equation}
The upper inequality can be summed over any set of child values.
The lower inequality is asserted only for retained children.
In particular, fix $0<\rho<1$.  If on $\mathcal H$ the original
conditional surprises for this pair lie in $[d_-,d_+]$ in units
$L=\log(1/\rho)$, then
every retained parent has at most $\rho^{-d_+}$ retained children,
and its retained conditional child masses obey
\[
 \delta\rho^{d_+}
 \le\pi(B=b\mid A=a)\le\delta^{-1}\rho^{d_-}.
\]
When a parent is nested only up to bounded ambiguity, include it in
the child symbol before applying this assertion.
\end{lemma}

\begin{proof}
Let
$\Lambda_{\mathrm{pair}}=\mu^0\otimes\lambda^0$ and
$\Lambda_{\mathrm{coord}}(\theta,q)
 =\lambda^0(\theta)\prod_j\xi_j^\theta(q_j)$.
These are probability laws on their respective observation spaces.
Use Lemma~\ref{lem:zero-information} with cutoff $v$ for both
likelihood ratios.  Pull both allowed sets back to the original
$(P,\theta)$ space through the corresponding observation maps and
intersect them with $\mathcal H$.  The resulting unnormalized
subprobability $r^0$ has mass at least $\underline M$, satisfies
$r^0\le p^0$, and obeys
\[
 r^0_{P,\theta}\le e^v\mu^0\lambda^0,\qquad
 r^0_{\theta,Q_1,\ldots,Q_k}
       \le e^v\lambda^0(\theta)\prod_j\xi_j^\theta(Q_j).
\]
The pullback is important: the restrictions for both comparisons
act on the same pair incidences, including when several points have
the same coordinate tuple.

Starting with $r^0$, repeatedly delete all incidences with $F=f$
whenever its current positive marginal is smaller than
$\delta p^0_F(f)$, for any $F\in\mathcal F$.  Each indexed value is
deleted at most once.  Charging its removed mass to $(F,f)$ bounds
the total loss by
$\delta\sum_{F\in\mathcal F}\sum_f p^0_F(f)=N\delta$.
Thus the final subprobability $r$ has mass at least
$\underline M-N\delta=\underline M/2$ and satisfies
\eqref{eq:block-core-floors}.  Restriction preserves both likelihood
upper bounds and $r\le p^0$.

For the pair product, the floors give
$\mu^0\le M\mu/\delta$ and
$\lambda^0\le M\lambda/\delta$ on the retained values.
Substitute them in the first likelihood upper bound and divide by
$M$ to obtain \eqref{eq:block-core-pair}.
For a surviving label and each coordinate value in its joint support,
\[
 \xi_j^\theta(Q_j)
 \le \frac{r_{\theta,Q_j}(\theta,Q_j)}
          {\delta\lambda^0(\theta)}
 =\frac{r_\theta(\theta)}{\delta\lambda^0(\theta)}
       \pi(Q_j\mid\theta).
\]
Conditioning the second likelihood upper bound on this label gives
\[
 \pi(Q_1,\ldots,Q_k\mid\theta)
 \le e^v\delta^{-k}
       \left(\frac{r_\theta(\theta)}{\lambda^0(\theta)}\right)^{k-1}
       \prod_j\pi(Q_j\mid\theta).
\]
Since $r\le p^0$, the parent ratio is at most one.  This proves
\eqref{eq:block-core-coordinate}; in particular the repeated label
in the factors has not been counted as an independent copy.

Finally, for a retained child,
\[
 \frac{\pi(B=b\mid A=a)}{p^0(B=b\mid A=a)}
 =\frac{r_B(b)/p^0_B(b)}{r_A(a)/p^0_A(a)}.
\]
Both ratios on the right lie in $[\delta,1]$, proving
\eqref{eq:block-core-prefix-ratio}.  Every retained child is
attained in $\mathcal H$, where its original conditional mass is
at least $\rho^{d_+}$ and at most $\rho^{d_-}$.  Summing the lower
original masses proves the support count, and substitution proves
the final bounds.
\end{proof}

Apply Lemma~\ref{lem:affine-block-core} to $p_n^0$ and $\mathcal H_n$.
Its list $\mathcal F_n$ contains the constant observation, $P$,
$\theta$, every $(\theta,Q_{j,\theta})$ at full block precision,
and every tested canonical prefix and joint numerator or denominator
symbol.  The parameter children use the fixed dyadic tree; every other
child includes its parent when needed, so the nested-symbol conclusion
applies exactly.  Write $N_n=|\mathcal F_n|$.
Set
\[
 t_n=\sqrt{d_n+L_n^{-1}},\qquad v_n=t_nL_n,\qquad
 M_{*,n}=1-\zeta_n-2t_n,\qquad
 \delta_n=\frac{M_{*,n}}{2N_n}.
\]
For all sufficiently large $n$, $M_{*,n}>0$ and
\[
 \frac{D_{\mathrm{pair}}+D_{\mathrm{coord}}+2/e}{v_n}
 \le 2t_n.
\]
The lemma therefore gives a restriction $r_n\le p_n^0$ of mass
$M_n\ge M_{*,n}/2$, supported on the simultaneous prefix and
likelihood trims.  Define its normalized pair law and actual marginals
\[
 \pi_n=r_n/M_n,\qquad \mu_n=(\pi_n)_P,\qquad
 \lambda_n=(\pi_n)_\theta,\qquad
 \xi_{j,n}^{\theta}=(\pi_n)_{Q_{j,\theta}\mid\theta}.
\]
The deletion proof acts on the common pair space; thus a retained
pair is an incidence for the canonical shear, and its support can
be used when labels are sampled again.

The finite lists can be increased slowly enough that
$\log N_n=o(L_n)$, since the precision is chosen after each list.
The one deterministic envelope
\[
 \kappa_n=
 \omega_n+t_n+
 \max\{2,|\mathcal I|\}\frac{\log(2N_n/M_{*,n})}{L_n}
 +\frac{\log(2/M_{*,n})}{L_n}
 \longrightarrow0
\]
bounds every normalized loss just obtained.  Suppress $n$ and put
$K_0=\rho^{-\kappa_n}$ in the following formulas.  The output is
the single finite law
\begin{equation}
 \begin{aligned}
 \pi(P,\theta)&\le K_0\mu(P)\lambda(\theta),\\
 \pi\bigl((Q_{j,\theta})_{j\in\mathcal I}\mid\theta\bigr)
       &\le K_0\prod_{j\in\mathcal I}\xi_j^\theta(Q_{j,\theta}).
 \end{aligned}
 \label{eq:rate-one-label-product}
\end{equation}
where every factor is an actual marginal of $\pi$, and the second
bound holds at every retained full label.  It also holds after
pushing any coordinate to a tested prefix.

For every retained value of each tested nested parent and child,
the ratio between its conditional mass in $\pi$ and in $p^0$
lies in $[\delta_n,\delta_n^{-1}]$.  In particular the parameter
and coordinate prefix bounds include
\begin{equation}
 \begin{aligned}
 \lambda(c_d\in I)&\le\rho^{sd-\kappa_n},&
 \lambda(r_d\in I)&\le\rho^{\sigma d-\kappa_n},\\
 \lambda(\Theta_b\mid\Theta_a)
       &\le\rho^{s(b-a)-\kappa_n}\quad(m=1),&
 \xi_j^\theta((Q_{j,\theta})_d=q)
       &\le\rho^{a_jd-\kappa_n}.
 \end{aligned}
 \label{eq:rate-parameter-bounds}
\end{equation}
Here the bounds refer to the finite tested list; $I$ is an interval
bin at the stated depth.  The first two inequalities are absolute
parameter masses; the conditional parameter and coordinate inequalities
are on retained parent fibers.  The corresponding original lower masses
on $\mathcal H_n$ give
\begin{equation}
 \#\{q:(Q_{j,\theta})_d=q
             \text{ for some retained neighbor of }\theta\}
 \le\rho^{-a_jd-\omega_n}
 \label{eq:rate-coordinate-support}
\end{equation}
for every retained full label and tested depth.  Analogous counts
hold for every listed joint child.  These support bounds are in
place before any label is replicated.  The upper mass bounds can
be summed over intervals or balls; a finite mesh of tested prefix
depths gives the bounds for intermediate radii with the mesh error.
The mesh is fixed before its precision is selected.

All later scalar and direction profiles are extracted from this
already chosen sequence $\pi_n$ at $\rho_n\to0$.  They do not
choose new $h_n$.  If a later profile selects a real parent depth
and a positive derivative increment, extend the nested surprises
from the rational prefixes and impose a fresh uniform event for
that finite real list before the restrictions which use it.
For each fixed increment, errors tending to zero in block units
also tend to zero after division by that increment.  Only then may
the increment shrink.  This is a subsequence of the existing block
sequence, not a new trace-scale choice.

An arbitrary later restriction does not automatically preserve
products of its new actual marginals.  If
$\widehat\pi=w\pi/m_R$ with $0\le w\le1$, then
\eqref{eq:rate-one-label-product} bounds its pair and coordinate
laws by $K_0/m_R$ times the old named products.  Decomposing the
two relative entropies against those products shows
\[
 \Mutn^{\widehat\pi}(P;\theta)\le\log(K_0/m_R),\qquad
 \sum_j\Entn^{\widehat\pi}(Q_j\mid\theta)
       -\Entn^{\widehat\pi}((Q_j)_j\mid\theta)
       \le\log(K_0/m_R).
\]
Thus a subpower restriction retains small information, while its
own product bounds are obtained by applying
Lemma~\ref{lem:affine-block-core} again after the required finite
prefix refresh.  The old upper prefix masses are first preserved
on typical parent fibers by Lemma~\ref{lem:typical-conditioning};
support counts persist by inclusion.  This is the procedure used
when the tied-clock argument later selects a homogeneous $Y$ profile.
In every later use of this one-label recoring procedure, reset
$\pi,\mu,\lambda$ and $\xi_j^\theta$ to its new law and actual factors.
Also renew $K_0$ and one common deterministic finite-test subpower
envelope for the products and tested prefix losses in the current
working units.  A fixed positive shortening by a factor $a$ preserves
subpower loss, since $o(L)=o(aL)$ when both are expressed in the
shortened units.

\subsection{Replication and the first projection tests}

For $1\le k_{\mathrm{rep}}\le3$, sample labels
$\theta_0,\ldots,\theta_{k_{\mathrm{rep}}-1}$ independently given
$P$ with the conditional kernel of $\pi$.  Their exact star law is
\[
 \Pi(P,\theta_0,\ldots,\theta_{k_{\mathrm{rep}}-1})
 =\mu(P)\prod_{j=0}^{k_{\mathrm{rep}}-1}\pi(\theta_j\mid P).
\]
Since $\pi(\theta\mid P)\le K_0\lambda(\theta)$, it satisfies
\begin{equation}
 \begin{aligned}
 \Pi&\le K_0^{k_{\mathrm{rep}}}\mu\prod_j\lambda_j,\\
 \Pi&\le K_0^{k_{\mathrm{rep}}-1}
       \pi(P,\theta_i)\prod_{j\ne i}\lambda_j
       \quad(0\le i<k_{\mathrm{rep}}).
 \end{aligned}
 \label{eq:rate-star-comparisons}
\end{equation}
where every $\lambda_j$ is a copy of $\lambda$.
These are the named comparison laws used after replication.

Let $R=w\Pi/m_R$ be a graph restriction chosen before conditioning.
Fix a set $J$ of full labels and leave $J^c$ free.  The conditioning
calculation in Theorem~\ref{thm:finite-conditioning}, applied to
\eqref{eq:rate-star-comparisons}, has two consequences on typical
fixed-label values.  First, for each $i\in J$, the point law is
dominated by a subpower times $\pi(P\mid\theta_i)$; pushing this
comparison through the coordinate map at label $i$ retains its
product against the named factors $\xi_j^{\theta_i}$.
Second, after the theorem's point-likelihood trim, the joint law
of the point and the free labels is compared with its actual
conditional point marginal times the copies $\lambda_j$ for
$j\notin J$.

Feed all coordinate products needed at these fixed labels and the
point--free-label comparison into one application of
Lemma~\ref{lem:marginal-core}.  Repeated occurrences of a coordinate
with different named factors are counted separately: for example,
$\xi_X^{\theta_0}$ and $\xi_X^{\theta_1}$ may differ even though
the map $X_\theta=X$ is invariant.  The simultaneous core replaces
both by the same current actual point marginal of $X$.  If
$J=\{0,1\}$ in the tied case, the resulting point law $\nu$ satisfies
\begin{equation}
 \nu_{X,U_0}\le K_1\nu_X\nu_{U_0},\qquad
 \nu_{X,U_1,Z_1}\le K_1\nu_X\nu_{U_1}\nu_{Z_1},
 \label{eq:rate-both-coordinate-products}
\end{equation}
and the full coordinate product at either label is available when
required.  Its coordinate marginal prefix bounds retain the rates
$x,q,z$ with subpower losses.  If label $2$ is free, the same core
gives a point--label-2 comparison with the actual point marginal
and a retained label comparison whose absolute prefix masses are bounded
by a subpower times those of the named pre-core label law.  Conditional
child caps for the new label marginal also require the tested parent
floors or typical-parent likelihood thresholds; these are imposed for
the finite list before those caps are used.

Here and in later uses, the number of replicas is fixed and the
finite factor list, graph retention $m_R$, likelihood thresholds,
and core retention are chosen with common deterministic subpower
envelopes.  The constants in Theorem~\ref{thm:finite-conditioning}
are fixed products of $K_0$, their inverse retentions and
thresholds, and the finite factor count.  Their normalized
logarithms therefore tend to zero together on every retained
conditional cell.  The bound does not depend on the number of
possible fixed-label values.
For example, if
\[
 \log(1/m_R),\ \log(1/\eta),\ \log(1/u),\
 \log(2N_{\mathrm{fac}})\le b_nL_n
 \qquad\text{with }b_n\to0
\]
throughout a chosen conditional array, then its fixed-label
comparisons and fixed-size cores have constants at most
$\exp(C(\kappa_n+b_n)L_n)$, with $C$ depending only on the
fixed numbers of replicas and factors.  This supplies one common
subpower envelope for those cells.

A subsequent point observation is fixed only if it is a function
of $P$ and the labels already in $J$.  Because the point factor
after the core is actual, Corollary~\ref{cor:point-slicing}
preserves the same pre-slice free-label comparison on every
positive point slice.  Any coordinate product needed within such
a slice is separately conditioned by
Lemma~\ref{lem:typical-conditioning}, on typical values of the
specified coordinate parent.  An observation involving a free
label must instead be included among the fixed labels before this
argument, or be handled by a different verified comparison.

Suppose first that $m\ne1$.  Use two replicas and fix a typical
label $0$, leaving label $1$ free.  The fixed-label core just
described gives an actual point law for $(X,Y_0)$ with ball
dimension $x+y$: its coordinate product and the tested prefix
upper masses bound a square of radius $\rho^d$ by
$\rho^{(x+y)d-o(1)}$.  The free comparison is the retained
pre-slice label-$1$ law, whose time marginal has dimension $s$.
The dominated incidence law has graph product mass
$\rho^{o(1)}$.  For each full label $1$ the projected coordinate
\[
 Y_1=Y_0+(c_1-c_0)X
\]
has at most $\rho^{-y-o(1)}$ bins on its retained neighbors by
\eqref{eq:rate-coordinate-support}.  The other component of the
label is kept as a tag.  Theorem~\ref{thm:planar-projection} gives
\begin{equation}
 y\ge\min\{1,x+y,(x+y+s)/2\}.
 \label{eq:rate-first-chain}
\end{equation}

For the chain from $w$ to $z$, fix also the full block symbols
$X,Y_0$.  The label-$0$ coordinate product, conditioned on
typical such values, gives the point population $(W_0,Z_0)$
dimension $w+z$.  Corollary~\ref{cor:point-slicing} keeps the
same free label-$1$ comparison.  The last coordinate is
$Z_0+(c_1-c_0)W_0$ plus a translation determined by the fixed
symbols and the full free label, and its original support count
is still available.  The projection theorem gives
$z\ge\min\{1,w+z,(w+z+s)/2\}$.
Interchanging the two shears gives the chains from $x$ to $w$
and from $y$ to $z$ with $\sigma$ in place of $s$; for the latter
the fixed point symbols are $X,W_0$.  All slices are finite
grid symbols, and the typical-conditioning loss has no factor
depending on their alphabet sizes.

For any of these inequalities, if the source $d$ is positive and
the target $e$ is below $1$, the alternative $e\ge d+e$ is
impossible.  The remaining alternative gives $e\ge d+s$ or
$e\ge d+\sigma$, as required.  This proves the distinct-clock
part of Proposition~\ref{prop:rate-transfer}.

Now let $m=1$. Equation~\eqref{eq:rate-shear} becomes
\begin{equation}\label{eq:rate-tied-shear}
 U_1=U_0+(\theta_1-\theta_0)X,\qquad
 Z_1=Z_0+\operatorname{flip}(\theta_1-\theta_0)\cdot U_0
                +\operatorname{prod}(\theta_1-\theta_0)X,
\end{equation}
where $\operatorname{flip}(c,r)=(r,c)$ and
$\operatorname{prod}(c,r)=cr$.
The trace controls the joint block $U$, so it supplies no separate
affine rates for $Y,W$.  Homogenize the conditional $Y$-prefix
probabilities on a subpower part of the current one-label law, and
write $g(d)$ for their limiting entropy in block units.  This is
a selection within the already fixed block sequence.  At each finite
stage take the post-split law as the new $p^0$.  Relative to the old
one-label law, select typical values of every tested parent needed
for the parameter and $X$ upper masses.  For each tested nested
parameter pair $\Theta_a,\Theta_b$, also impose the capacity tail
\[
 p^0(\Theta_b\mid\Theta_a)
       \ge\tau C^{-1}\rho^{2(b-a)}.
\]
Here $C=4$ is the capacity constant of the fixed parameter tree.
There are at most $C\rho^{-2(b-a)}$ grid children over each parent,
so the failure of this condition has $p^0$-mass at most $\tau$.
Choose $\tau$ with $N_{\mathrm{test}}\tau\to0$ and
$\log(1/\tau)=o(L)$, where $N_{\mathrm{test}}$ is the number of these
tested pairs and $L=\log(1/\rho)$.  Intersect all good-parent filters,
capacity tails, and fresh $Y$ and coordinate surprise bins to form
one allowed event $\mathcal H$ for $p^0$, and keep
$\boldsymbol1_{\mathcal H}p^0$ unnormalized until the repeated
application of Lemma~\ref{lem:affine-block-core}.  Its list contains
all the tested numerator and denominator observations.  The old
named products give the small information needed for this extraction;
the simultaneous floors against $p^0$ preserve the inherited upper
masses and selected $Y$ exponents on the same retained fibers.
The new products use the actual marginals of the output, and its
support is a subset of the original coordinate support.

Relabel this final one-label law as $\pi$, with actual point and
parameter marginals $\mu,\lambda$.  If $\delta$ is the repeated core's
threshold, its simultaneous prefix floors give, on every retained
tested parameter child,
\begin{equation}
 \lambda(\Theta_b\mid\Theta_a)
 \ge\delta\tau C^{-1}\rho^{2(b-a)}
 =\rho^{2(b-a)+o(1)}.
 \label{eq:rate-parameter-capacity-floor}
\end{equation}
This uses only the planar child capacity and supplies a lower mass
for the short mesh denominators used below.  Put
\[
 g_*=\inf_{0<d\le1}\frac{g(d)}d.
\]
The law $\pi$ has the $X$--$U$ product, the $X$ and time upper
masses, and the homogeneous $Y$ prefix masses and support counts
on the same retained incidences.  Each $(P,\theta_i)$ marginal of
its exact star is this same law.  If a later graph restriction is
$R=w\Pi/m_R$, that marginal is bounded by $\pi/m_R$.
Typical label likelihood thresholds therefore preserve its named
prefix upper bounds, and support counts can only decrease.
Pushing the $X$--$U$ product through $U\mapsto Y$ bounds a square
of radius $\rho^r$ in $(X,Y_0)$ by
$\rho^{xr+g(r)-o(1)}$ on the retained typical labels.
Both replicas use the profile of this one law.

At precision $\rho^d$, the $(X,Y_0)$ point population has ball
dimension at least $x+g_*$, whereas its $Y_1$ projection count has
exponent $g(d)/d$. Indeed $g(r)\ge g_*r$ at every smaller prefix;
the requisite ball bounds follow from the homogeneous conditional
masses. Fix label $0$, leave label $1$ free, and use the same product
comparison as for \eqref{eq:rate-first-chain}. At each fixed $d>0$,
Theorem~\ref{thm:planar-projection} gives
\[
 \frac{g(d)}d\ge\min\{1,x+g_*,(x+g_*+s)/2\}.
\]
Taking the infimum does not assume it is attained: test successively
fixed $d$ with $g(d)/d$ approaching the infimum, choosing all losses
after $d$. If $x>0$, the resulting inequality implies
$g_*\ge\min(1,x+s)\ge s$. Since $g(1)\le q$, we have proved
$q\ge s$.

\subsection{Directions between parameter samples}

The reciprocal use of pinned directions and the local linearization
below have close precedents in the radial-projection bootstrap of
Shmerkin and Wang and the thin-tube bootstrap of Orponen, Shmerkin,
and Wang \citep{ShmerkinWang2025Distance,OrponenShmerkinWang2024Radial}.
We prove the version needed here for the retained conditional pair law,
including its selection order and the resulting improvement of the
directional rate.

Throughout this subsection and its applications, set $\dir(0)=(1,0)$
and give this fixed unit vector the ordinary nested direction-grid
prefixes.  Thus the direction and scalar observations are defined on
the entire finite space, including coincident labels.

We give the needed directional statement with its assumptions in
finite probability form. A child bound of rate $s$ means the conditional parameter-prefix
bound for $\lambda(\Theta_b\mid\Theta_a)$ in
\eqref{eq:rate-parameter-bounds} at every fixed tested pair of depths,
on uniform trims and typical parent fibers.  On each tested good parent,
it applies to the entire named conditional population, hence to all its
positive child values.  This includes
all intermediate radii after using a slowly growing finite mesh.

\begin{proposition}[Directional bootstrap]\label{prop:directional-bootstrap}
Let $0<s\le1$. Let a sequence of pair laws of bounded planar
parameters $(A,B)$ be dominated, with subpower loss at precision
$\rho$, by the product of two populations having conditional child
bounds of rate $s$.  Assume that for each fixed $t>0$, both comparison
marginals satisfy
\begin{equation}\label{eq:rate-no-strip}
 \sup_{\ell}\mu\{\dist(\cdot,\ell)\le\rho^t\}
       \le\rho^{\kappa_t+o(1)}\quad\text{for some }\kappa_t>0,
\end{equation}
and likewise for the other marginal. Suppose the direction-prefix
probabilities conditional on either pin are homogenized on the same
pair law.  For each finite test, assume a common event of subpower pair
mass on which the required original direction bins are uniform and the
pair support lies over the simultaneous tested good-parent fibers for
the child bounds.  Then, on uniform trims, the direction of $B-A$ conditional
on either pin has prefix ball dimension at least $s$.
The assertion is hereditary under subpower restrictions which retain
the stated child bounds and product comparisons.
\end{proposition}

\begin{proof}
Let $\mathsf P$ be the already homogeneous pair law in the hypothesis,
with $\mathsf P\le K_{\mathrm{pair}}\mu_A\otimes\mu_B$.
Write $F_B(d)$ for its conditional direction entropy given the full
pin $B$, and similarly $F_A(d)$.  The nested direction grids make both
profiles Lipschitz, vanishing at zero and absolutely continuous.
For each finite test choose one event $\mathcal H$ of mass
$m_H=\mathsf P(\mathcal H)>0$ carrying both the original conditional
direction bins and the simultaneous good-parent filters from the
hypothesis.  Keep
$\mathsf P,\mathcal H$ named, and put
$\mathsf R=\boldsymbol1_{\mathcal H}\mathsf P/m_H$.
The inverse retention and all later thresholds use the common
subpower envelope for this test.

Fix a differentiability point $j\in(0,1)$ of $F_B$ and a much smaller
fixed increment $a>0$.  For $0<\eta<1/2$, apply
Lemma~\ref{lem:typical-conditioning} with the observation $A_j$, both
to $\mathsf R\le(K_{\mathrm{pair}}/m_H)\mu_A\otimes\mu_B$ and to its
retention from $\mathsf P$.  Outside $\mathsf R_{A_j}$-mass at most
$2\eta$, the conditional law $\mathsf R_j=\mathsf R(\,\cdot\mid A_j)$
and $r_j=\mathsf P(\mathcal H\mid A_j)$ satisfy
\[
 \mathsf R_j\le K_j\mu_A(\,\cdot\mid A_j)\otimes\mu_B,
 \qquad K_j=\frac{K_{\mathrm{pair}}}{m_H\eta},\qquad r_j\ge m_H\eta.
\]
There is no inverse original parent mass.  Fix such a parent among the
tested child-bound fibers supplied by the hypothesis, with center
$a_0$, and write $A=a_0+\rho^j\xi$.  Its named rescaled child law has
ball dimension at least $s$ at precision $\rho^a$.  Restrict to
$|B-a_0|\ge\rho^{o(1)}$ using the displayed product and the
positive-dimensional ball bounds at each fixed positive depth.
If its relative retained mass is $v_j$, write $\mathsf G_j$ for the
normalized restricted graph; then
\[
 \mathsf G_j\le\frac{K_j}{v_j}\mu_A(\,\cdot\mid A_j)\otimes\mu_B,
\]
with $v_j^{-1}$ subpower.

For each pin, the number of retained depth-$(j+a)$ direction bins
inside a depth-$j$ bin is at most
\begin{equation}\label{eq:rate-direction-count}
 \rho^{-[F_B(j+a)-F_B(j)]-o(1)}.
\end{equation}
This follows by dividing the original $\mathsf P$ parent and child
masses conditional on that pin, using their uniform bins on
$\mathcal H$.  Further fixing $A_j$
only removes support. Computing the direction from $a_0$ instead of
$A$ determines the coarse direction bin to subpower ambiguity.
Taylor expansion of the direction map shows that, in child units,
the finer bins give intervals for the scalar projection of $\xi$
perpendicular to $B-a_0$, with error
\[
 O(\rho^a)+\rho^{j-o(1)}.
\]
Indeed the second derivative of the direction map is
$O(|B-a_0|^{-2})$; dividing its $O(\rho^{2j}|B-a_0|^{-2})$
error by the first-order scale $\rho^j/|B-a_0|$ gives the second
term. It is smaller than the working precision when $a\ll j$.
All translations and dilations have subpower size.

Before resampling, apply Lemma~\ref{lem:marginal-core} to the two named
factors of $\mathsf G_j$.  With threshold $1/4$ it gives a normalized
pair $q_j$ of raw mass $M_j\ge1/2$ satisfying
\[
 q_j\le16\frac{K_j}{v_j}(q_j)_A\otimes(q_j)_B.
\]
Its actual child and pin marginals are bounded respectively by
$(K_j/(v_jM_j))\mu_A(\,\cdot\mid A_j)$ and
$(K_j/(v_jM_j))\mu_B$, so their ball and strip bounds persist; its
support retains the direction counts.  Now sample two pins independently
over $(q_j)_A$ using the kernel of $q_j$.
Theorem~\ref{thm:finite-conditioning} compares the star with the
independent product of these actual marginals
and with either retained pair $q_j$ times the other pin marginal.
For any fixed $\tau>0$, an absolute sine of the angle at most $\rho^{\tau a}$ between
the two directions puts the second pin in a strip through $a_0$ of
width $O(\rho^{\tau a})$. Equation~\eqref{eq:rate-no-strip} makes
this event have probability tending to zero, even after the indicated
product losses. Angles are thus tested modulo $\pi$, including
nearly antiparallel directions. Fix a typical pair with the absolute
sine of its angle at least $\rho^{\tau a}$. The child
retains its ball bound, and its two projections have counts
\eqref{eq:rate-direction-count}. Two such projection intervals locate
it within a ball of radius $O(\rho^{(1-\tau)a})$. Therefore
\[
 2\,[F_B(j+a)-F_B(j)]\ge s(1-\tau)a-o(1).
\]
Here $j,a,\tau$ are fixed before $\rho\to0$. Then let $a\downarrow0$
and $\tau\downarrow0$. We obtain $F_B'(j)\ge s/2$ almost everywhere.
The same proof for the opposite pin uses the \emph{same} homogeneous
pair law. Integration gives $F_A(d),F_B(d)\ge sd/2$.

Suppose now that both profiles have the lower bound $ud$, where
$0<u\le s$.  The original conditional bins on $\mathcal H$ give the
upper bound $\rho^{ur-o(1)}$ for each retained direction bin of depth
$r$ given its pin, at every fixed required prefix.
For this reciprocal step use $\mathsf G_j$ before the auxiliary
two-pin core.  Since $A_j$ is determined by $A$, the conditional law
of $B$ given $A$ under $\mathsf P(\,\cdot\mid A_j)$ is still that under
$\mathsf P$.  At depths $r\le a\ll j$, the directions about $A$ and
$a_0$ differ by less than the tested precision on the separation set.
Average the original conditional upper bound on $\mathcal H$ over
$A$ in this parent before normalizing either restriction.  The actual
pin marginal $(\mathsf G_j)_B$ therefore obeys the direction bound
$\rho^{ur-o(1)}$ with the extra factor $(r_jv_j)^{-1}$, which is
subpower.  Apply Lemma~\ref{lem:marginal-core} directly to
$\mathsf G_j$ with its two named factors, again with threshold $1/4$.
If this core has raw mass $M\ge1/2$, its actual pin marginal is at most
$M^{-1}(\mathsf G_j)_B$, while its actual child marginal is at most
$(K_j/(v_jM))\mu_A(\,\cdot\mid A_j)$.  The former keeps the reciprocal
direction bound and the latter the child ball bound; the core supplies
their actual product comparison.  Thus
Theorem~\ref{thm:planar-projection} applies with dimensions $s,u$ and
count \eqref{eq:rate-direction-count}.  It gives
\[
 F_B'(j)\ge\min\{1,s,(s+u)/2\}=(s+u)/2,
\]
and the same estimate for $F_A$. Integrate and repeat. Starting at
$u=s/2$, the finite iterates increase to $s$. For each prescribed
error only finitely many repetitions are needed. All tested parents,
increments, radii, and repetitions are fixed before choosing the
small-parameter losses. A diagonal over those finite choices gives
$F_A(d),F_B(d)\ge sd$.  Apply
Lemma~\ref{lem:typical-conditioning} separately at each full pin to
the common original bin event $\mathcal H$.  On the same normalized
law $\mathsf R=\mathsf P(\,\cdot\mid\mathcal H)$, outside pin mass
at most $\eta$ for each choice of pin, the actual conditional direction
caps hold with the additional factor $(m_H\eta)^{-1}$.  This proves
the claimed conditional mass bounds without normalizing an intersection
of the two pin filters.
\end{proof}

\subsection{Three labels and the full-precision slices}

Assume first that \eqref{eq:rate-no-strip} holds for the one-label
parameter population.  Use the globally defined parameter-tree atoms
from the one-label construction.
Their simultaneous bins, parent filters, and prefix floors give the
rate-$s$ upper child cap in \eqref{eq:rate-parameter-bounds} and the
lower capacity bound \eqref{eq:rate-parameter-capacity-floor} for the
same current $\lambda$ on every retained tested child.  These are the
conditional masses of the named one-label population used below.

They cover a later fixed real parent $j$ and increment $a$ by a finite
mesh.  Choose tested depths $j_-<j<j_+<b_-<j+a$ within mesh size
$\delta_{\mathrm{mesh}}$ of $j$ and $j+a$.  For the compatible nested
prefixes of a retained child, use the globally defined tree atoms at
all five depths below.  In particular the denominator uses
the already floored $\Theta_{j_+}$ atom, unchanged by the choice of the
real depth $j$.  Thus
\[
 \lambda(\Theta_{j+a}\mid\Theta_j)
 \le\frac{\lambda(\Theta_{b_-}\mid\Theta_{j_-})}
          {\lambda(\Theta_{j_+}\mid\Theta_{j_-})}
 \le\rho^{s(b_--j_-)-2(j_+-j_-)-o(1)}.
\]
Choose $\delta_{\mathrm{mesh}}\ll a$ after $j,a$ are fixed, and then
the precision.  This gives the required rate-$s$ child cap with the
mesh error and one common envelope, within the existing block sequence.

Sample three labels over $P$ using \eqref{eq:rate-star-comparisons}.
At each finite stage apply Lemma~\ref{lem:homogenization} to the
direction-grid children of the first pair conditional on each full pin,
and to the scalar-grid children below conditional on the full pair.
Their child alphabets have size at most a fixed power of $\rho^{-1}$
for each finite list, while no bound on the parent alphabets is needed.
Choose the lemma's parent and child thresholds
together for this finite list, with vanishing total exceptional mass
under the common envelope.  The diagonal in
Lemma~\ref{lem:profile-diagonal} gives a subpower retained triple law
carrying the actual conditional profiles.
Its first-pair marginal is dominated by a subpower times the unchanged
named product $\lambda_0\lambda_1$.  These copies retain the tested
child and no-strip bounds just verified, and the two direction profiles
belong to this same first-pair law.  Thus
Proposition~\ref{prop:directional-bootstrap} applies.  Lift its common
pair event to this same point--label law and normalize it as
$\mathcal R$.  Its pair marginal is the normalized law in the
proposition, so the actual conditional direction caps hold on the
respective good-pin values.  Keep this pre-average law for the two
projection tests below.  Write
\[
 d=\dir(\theta_1-\theta_0),\qquad
 \zeta_2=\operatorname{flip}(d)\cdot\theta_2.
\]
For the scalar test, retain one good label-$0$ set from the proposition,
of $\mathcal R$-mass at least $1-\eta_{\mathrm{pin}}$.
The label marginal satisfies
$\mathcal R_{012}\le K_{\mathrm{lab}}\lambda_0\lambda_1\lambda_2$
with subpower $K_{\mathrm{lab}}$.  Also retain the pair values
\[
 \mathcal R_{01}(b)\ge
       \eta_{\mathrm{pair}}(\lambda_0\lambda_1)(b).
\]
Their complement has $\mathcal R$-mass at most
$\eta_{\mathrm{pair}}$ by \eqref{eq:conditioning-likelihood}.
Normalize the intersection of these two pair-only filters as
$\mathcal R^0$.  Averaging the conditional direction cap over the
retained good label-$0$ values gives an absolute rate-$s$ cap for
$d$ under $(\mathcal R^0)_{01}$, with full-pair tags and only the
inverse retained mass loss.  Since both filters depend on the pair,
their normalization gives the exact comparison
\[
 (\mathcal R^0)_{012}
 \le\frac{K_{\mathrm{lab}}}{\eta_{\mathrm{pair}}}
       (\mathcal R^0)_{01}\lambda_2.
\]
The third population $\lambda_2$ has point dimension at least $s$.
The thresholds tend to zero with subpower inverses in the common
finite-test envelope.

Let $l(r)$ be the inherited scalar profile of $\zeta_{2,r}$ given the
full pair under $\mathcal R^0$.  The pair-only filters leave its
conditional law on each retained pair unchanged, and the earlier
common bins give its profile under this subpower restriction.
For a finite list of $r$, choose one event $\mathcal H_{\mathrm{sc}}$
of mass $m_{\mathrm{sc}}=\mathcal R^0(\mathcal H_{\mathrm{sc}})$ on which
\[
 \left|-L^{-1}\log\mathcal R^0(\zeta_{2,r}\mid b)-l(r)\right|
       \le\omega_{\mathrm{sc}},\qquad \omega_{\mathrm{sc}}\to0,
\]
and put $\mathcal R^{\mathrm{sc}}
=\mathcal R^0(\,\cdot\mid\mathcal H_{\mathrm{sc}})$.
The inverse mass is subpower.  This law is dominated by
$K_{\mathrm{lab}}/(\eta_{\mathrm{pair}}m_{\mathrm{sc}})$ times the
named product $(\mathcal R^0)_{01}\lambda_2$, so the pre-event pair
population supplies the direction cap.  The original lower scalar
masses on $\mathcal H_{\mathrm{sc}}$ give at most
$\rho^{-l(r)-\omega_{\mathrm{sc}}}$ retained scalar bins per full pair.
Theorem~\ref{thm:planar-projection}, with full-pair tags and dimensions
$s,s$, therefore gives $l(r)\ge sr$.

Put $r_b=\mathcal R^0(\mathcal H_{\mathrm{sc}}\mid b)$ and
$\lambda_b^0=(\mathcal R^{\mathrm{sc}})_{\theta_2\mid b}$.
Lemma~\ref{lem:typical-conditioning} gives
$r_b\ge m_{\mathrm{sc}}\eta_{\mathrm{sc}}$ outside
$(\mathcal R^{\mathrm{sc}})_{01}$-mass at most $\eta_{\mathrm{sc}}$.
For each retained scalar bin on such a pair,
\[
 \lambda_b^0\{\zeta_{2,r}\in I\}
 \le r_b^{-1}\mathcal R^0(\zeta_{2,r}\in I\mid b)
 \le(m_{\mathrm{sc}}\eta_{\mathrm{sc}})^{-1}
       \rho^{l(r)-\omega_{\mathrm{sc}}}
 \le\rho^{sr-o(1)}.
\]
Choose $\eta_{\mathrm{sc}}$ with a subpower inverse and intersect these
pairs with the typical fixed-pair values from the star comparisons for
the graph $\mathcal R^{\mathrm{sc}}$.  Also impose the pair-only filter
$|\theta_1-\theta_0|\ge\rho^{\vartheta_n}$, with
$\vartheta_n\downarrow0$ chosen slowly for this graph's envelope.
Indeed, for each fixed $0<t\le1$, its pair domination by a subpower times
$\lambda_0\lambda_1$ and the absolute parameter ball cap give
\[
 (\mathcal R^{\mathrm{sc}})_{01}
       \{|\theta_1-\theta_0|<\rho^t\}
 \le\rho^{st-o(1)}.
\]
A slow diagonal makes the excluded mass vanish at $t=\vartheta_n$.
This pair-only filter leaves conditional laws at retained pairs unchanged.
With $\alpha=|\theta_1-\theta_0|$, it gives
$\alpha^{-1}\le\rho^{-\vartheta_n}=\rho^{-o(1)}$.
On each remaining pair $b$, feed the
point--third-parameter comparison and the two coordinate products
supplied by the $i=0,1$ star comparisons into one simultaneous core.
Let $M_b$ be the total raw retained mass relative to
$\mathcal R^{\mathrm{sc}}(\,\cdot\mid b)$, including the point-likelihood
trim and all core deletions.  Then the new third marginal satisfies
$\lambda_{01}\le M_b^{-1}\lambda_b^0$.  All inverse retentions use the
same subpower envelope.  Hence
\begin{equation}\label{eq:rate-third-slope}
 \lambda_{01}\{\zeta_{2,r}\in I\}\le\rho^{sr-o(1)}.
\end{equation}
The two coordinate products now use actual marginals of that same
retained point law.  This establishes their coexistence with the scalar
slope bound before any point-coordinate slice.

Fix such a pair, and write $\theta_1-\theta_0=\alpha d$, where
$\alpha^{-1}$ has subpower size.  Before the last trim form the
real linear coordinates
\[
 N=d^\perp\cdot U_0=d^\perp\cdot U_1,\qquad
 L=d\cdot U_1.
\]
Bin $N,L$ at precision $\rho$.  Enlarge the $U_0$ symbol to
include the $N$ bin, which has boundedly many possibilities given
it, and use $(N,L)$ as the $U_1$ symbol.  Rotated and original
square bins determine one another with bounded ambiguity.

Let $\widetilde\nu^\star(P,\theta_2)$ denote the joint law just obtained
from the fixed-pair core, with actual point marginal $\widetilde\nu$.
It has
\[
 \widetilde\nu^\star(P,\theta_2)
 \le A\,\widetilde\nu(P)\lambda_{01}(\theta_2)
\]
with subpower $A$, and its point marginal carries both coordinate
products and the $X,Z_1$ prefix upper masses.  The bounded new
symbols change the coordinate multiinformations by $O(1)$ in
natural-log units.  Apply Lemma~\ref{lem:zero-information} and the
simultaneous marginal core to these two coordinate products on
the point law $\widetilde\nu$ alone.  The likelihood events and
all core deletions are functions of the point and the already fixed
pair.  They therefore give a point-only weight $0\le w(P)\le1$ of
subpower retained mass
$M=\mathbb E_{\widetilde\nu}w>0$.  Lift that weight to the joint
law and put
\[
 \nu^\star(P,\theta_2)=\frac{w(P)\widetilde\nu^\star(P,\theta_2)}M,
 \qquad
 \nu(P)=\frac{w(P)\widetilde\nu(P)}M.
\]
The point--free-label comparison is then exactly
\begin{equation}
 \nu^\star(P,\theta_2)
 \le A\,\nu(P)\lambda_{01}(\theta_2).
 \label{eq:rate-point-only-trim}
\end{equation}
There is no new factor $M^{-1}$ in this comparison: the same point
weight and normalization occur in its actual point marginal.
The coordinate core gives \eqref{eq:rate-both-coordinate-products}
for the new finite symbols on $\nu$, and the
$X,Z_1$ upper masses lose at most the subpower factor $M^{-1}$.

The fixed labels are still $J=\{0,1\}$ and label $2$ remains free.
Both $N$ and $X$ are finite functions of the point and these fixed
labels.  By Corollary~\ref{cor:point-slicing}, every positive
$(N,X)$ slice of \eqref{eq:rate-point-only-trim} keeps the same
pre-slice comparison $\lambda_{01}$ and the actual sliced point
marginal.  In particular its scalar slope bound remains available.
No exact real coordinate is conditioned on.

We prove carefully that the population $(L,Z_1)$ conditional on
$(N,X)$ has ball dimension at least $x+z$. If a $U_0$ cell with center
$u_0$ is fixed, then $L$ is within $O(\rho)$ of
$d\cdot u_0+\alpha X$. The first product in
\eqref{eq:rate-both-coordinate-products} and the marginal $X$ bound
give, for an interval $I$ of length $r\ge\rho$ (absorbing the
constant enlargement),
\begin{equation}\label{eq:rate-L-slice}
 \nu(L\in I\mid U_0=u_0)
        \le K_1 C\alpha^{-x}r^x,
\end{equation}
where $C$ is subpower. Average this inequality over $u_0$ conditional
on $N=n$. This proves the same bound for $L\mid N=n$.

The second product in \eqref{eq:rate-both-coordinate-products},
after rotating $U_1$ to $(N,L)$, gives the exact finite inequality
\begin{equation}\label{eq:rate-double-slice}
 \nu(l,z\mid N=n,X=t)
 \le K_1\frac{\nu_X(t)\nu_N(n)}{\nu_{X,N}(t,n)}
              \nu(l\mid N=n)\nu_{Z_1}(z).
\end{equation}
The $\nu$-probability of the values $(t,n)$ where the displayed
ratio exceeds $A>1$ is at most $A^{-1}$: sum
$\nu_{X,N}(t,n)<A^{-1}\nu_X(t)\nu_N(n)$ over those values.
On the remaining slices, \eqref{eq:rate-L-slice} and the $Z_1$
marginal bound imply
\begin{equation}\label{eq:rate-xz-mass}
 \nu((L,Z_1)\in I\times J\mid N=n,X=t)
       \le K_1^2 A C^2\alpha^{-x}r^{x+z}
\end{equation}
for intervals of length $r$. Take $A\to\infty$ with subpower
growth. Thus there is no loss proportional to the number of
full-precision $X$ or $N$ bins.

The free label-$2$ law on these slices uses the same comparison
$\lambda_{01}$ from \eqref{eq:rate-third-slope}: the point-only
identity \eqref{eq:rate-point-only-trim} and its $(N,X)$ slicing
were verified above.  Hence the free scalar slope bound is available
on these slices.

In \eqref{eq:rate-tied-shear}, the change from frame $1$ to frame
$2$ is a scalar projection of $(L,Z_1)$: its slope is
$\operatorname{flip}(d)\cdot(\theta_2-\theta_1)$, and the remaining
terms are translations determined by $(N,X)$ and the free label.
The support count of $Z_2$ per label is still at most
$\rho^{-z-o(1)}$, because every restriction only removed original
neighbors. Apply Theorem~\ref{thm:planar-projection} to
\eqref{eq:rate-xz-mass} and \eqref{eq:rate-third-slope}. We obtain
\begin{equation}\label{eq:rate-tied-xz}
 z\ge\min\{1,x+z,(x+z+s)/2\}.
\end{equation}
For $x>0$ and $z<1$ this is exactly $z\ge x+s$.

For the second test, return to a fresh copy of the pre-average
point--label law $\mathcal R$, restrict it to its good label-$0$ values
from Proposition~\ref{prop:directional-bootstrap}, and marginalize
label $2$.  This lifts the pair-only filter to the same incidence law;
the star comparisons and label-$0$ coordinate support remain on one
graph.  The conditional direction law of label $1$ at each retained
label $0$ is unchanged and has rate $s$.
Fix a typical label $0$ and apply the simultaneous core to the
point--label-$1$ comparison and the label-$0$ coordinate product.
Its new label-$1$ marginal retains the direction cap with a subpower
inverse mass loss.  Then fix the point coordinates $X,Z_0$; label $1$
is free.  The point population is $U_0$, of dimension $q$ on typical
slices by that coordinate product.  Its direction comparison is this
retained \emph{pre-slice} label-$1$ marginal, for the direction
$\operatorname{flip}(\theta_1-\theta_0)$.
The fixed-label core uses the $i=0$ comparison in
\eqref{eq:rate-star-comparisons};
Corollary~\ref{cor:point-slicing} preserves this same direction
population on the subsequent $(X,Z_0)$ slice.
Equation~\eqref{eq:rate-tied-shear} projects $U_0$ to the $Z_1$
support, giving
\begin{equation}\label{eq:rate-tied-q}
 z\ge\min\{1,q,(q+s)/2\}.
\end{equation}
If $q>2-s$, both other terms exceed or equal $1$, so $z\ge1$.
This proves the last assertion in the absence of strip concentration.

All rotations and scalar slices just used have a finite-grid meaning.
Rotated finest-scale squares meet boundedly many original squares
and conversely. Add the two descriptions and their bounded neighbor
tags when conditioning. Equation~\eqref{eq:rate-tied-shear} then
holds up to $O(\rho)$; all interval bounds use constant enlargements.
The factors $\alpha^{-1}$ and the harmless angular separations can
be bounded by $\rho^{-\eta}$ for each fixed $\eta>0$ before the
limit. Replacing a finest scale by a $\rho^{-O(\eta)}$ enlargement
changes the displayed exponent inequalities by $O(\eta)$, which is
sent to zero afterwards.

\subsection{Concentration in a strip}

It remains to justify the alternative omitted above. Test the original
retained one-label marginal $\lambda$ before sampling replicas. If
\eqref{eq:rate-no-strip} fails along some subsequence at a fixed
depth $t>0$, pass to a subsequence and choose a strip of width
$\rho^t$ whose $\lambda$-mass is $\rho^{o(1)}$. This follows
by choosing a sequence on which the nonnegative exponent of the
maximal strip mass tends to zero. Otherwise a positive lower bound
for that exponent supplies \eqref{eq:rate-no-strip} at each fixed
depth, and a diagonal supplies all the depths used above. Subpower
product domination then prevents concentration of a replicated
marginal that was absent from the one-label population.

Restrict the one-label graph to the chosen strip and normalize it as
the new $p^0$.  Its mass in the previous law is subpower, so the old
named products give the small information required by
Lemma~\ref{lem:affine-block-core}.  Shorten the working block to a fixed
positive depth no greater than $t$, if necessary, and fix its prefix
list before taking the small-parameter limit.  Intersect the fresh
bins with all typical-parent filters needed for the inherited
conditional bounds, keep this allowed restriction unnormalized, and
apply the lemma once for that finite list.  Relabel its output $\pi$
and its actual pair marginals $\mu,\lambda$.  Its pair and coordinate
products now use actual factors before any new replicas are sampled,
and its support still lies in the chosen strip.  In these shorter
block units the coordinate rates remain $x,q,z$, by
\eqref{eq:rate-post-affine}; the time and joint parameter rates remain
at least $s$.  The strip is now as thin as the finest precision.
Let $d_0$ be its unit direction. Both coordinates of $d_0$ have
exponent zero. Indeed, if (say) its first coordinate had magnitude
at most $\rho^b$ along a subsequence with $b>0$, the bounded strip
would lie in a time-coordinate interval of a fixed positive depth.
Its subpower retained mass contradicts the time marginal bound in
\eqref{eq:rate-parameter-bounds}. The second coordinate is treated
using the positive normal marginal bound. In particular
\[
 |\operatorname{flip}(d_0)\cdot d_0|
       =2|(d_0)_1(d_0)_2|=\rho^{o(1)}.
\]

The scalar position along the strip has prefix ball dimension at
least $s$: projection to the time coordinate has this bound, and
its slope along the strip has exponent zero. After sampling three
labels over the restricted graph, typical pairs have separation
$\rho^{o(1)}$ by this scalar bound. Their differences are parallel
to $d_0$ up to finest-scale errors, so the scalar third slope
$\operatorname{flip}(d_0)\cdot\theta_2$ has dimension at least $s$
directly. Typical fixed-pair comparisons preserve it. The proof of
\eqref{eq:rate-tied-xz}, including both products at labels $0,1$,
now applies with $d_0$ in place of $d$: perpendicular invariance and
parallel transfer have only finest-scale errors. Thus $x>0,z<1$
again gives $z\ge x+s$.

For the remaining test, start from a fresh copy of the strip star before
the preceding fixed-pair restrictions, then fix label $0$ and $X$,
leaving label $1$ free. Set $B_0=\operatorname{flip}(d_0)\cdot U_0$. If $q>2-s$,
then $q>1$. A strip of width $r$ for $B_0$ is covered by $O(r^{-1})$
balls of radius $r$ in the bounded $U_0$-plane. Hence its mass is
at most $C r^{q-1}$, with subpower $C$. The label-$0$ coordinate
product, after the typical $X$ slice, gives ball dimension at least
$q-1+z$ for $(B_0,Z_0)$. The free comparison law is the retained
label-$1$ marginal; its scalar position along the strip has dimension
$s$. Equation~\eqref{eq:rate-tied-shear} projects this pair to
$Z_1$ with that scalar slope, up to a known translation and
finest-scale errors. Theorem~\ref{thm:planar-projection} gives
\[
 z\ge\min\{1,q-1+z,(q-1+z+s)/2\}.
\]
If $z<1$ and $q>2-s$, each expression inside this minimum is
strictly greater than $z$: for the last one use
$q-1+s>1>z$. This is impossible. Consequently $z<1$ implies
$q\le2-s$ also in the strip case.

This completes the proof of Proposition~\ref{prop:rate-transfer}.
The order is the block sequence $(h_n,\eps_n)$ first, the
homogeneous scalar and direction profiles second, and their tested
depths and positive derivative increments third. At that third stage
every test fixes its finite list before taking a subsequence limit
in $n$, and uniform trims precede all restrictions using their
bounds. The derivative increment is sent to zero only afterwards.
Each prescribed accuracy in the directional bootstrap requires
finitely many iterations; a diagonal over those finite tests proves
the exact inequalities. Thus neither a new $h_n$ chosen from a
later direction profile nor a uniform differentiability modulus
across profiles is required.
